% Part: model-theory
% Chapter: interpolation
% Section: introduction

\documentclass[../../../include/open-logic-section]{subfiles}

\begin{document}

\olfileid{mod}{int}{int}

\olsection{مقدمة}

مؤدى مبرهنة الاستيفاء أنه، لجملتين، إذا كان
$\Entails !A \lif !B$، وجدت !!a{sentence}~$!C$ تحقق
$\Entails !A \lif !C$ و$\Entails !C \lif !B$،
مع شرط آخر: كل !!{constant} و!!{function} و!!{predicate}
غير~$\eq$ يرد في $!C$ يجب أن يرد في $!A$ و$!B$
جميعًا. وهذه الـ!!{sentence}~$!C$ تسمى \emph{مستوفية}
لـ$!A$ و$!B$.

وللمبرهنة شأن في نفسها، غير أن أعظم نفعها هنا ما يترتب
عليها من نتيجتين: إحداهما في قابلية التعريف داخل نظرية،
والأخرى في الشروط التي يجوز معها ضم نظريتين متسقتين
من غير أن يزول الاتساق. والأولى مبرهنة بث في قابلية
التعريف، والثانية مبرهنة روبنسون للاتساق المشترك.

\end{document}
